\documentclass{Beautybook-V6.1}
\usepackage{rotating}
\tikzset{>=Stealth}
\setlist{nosep,font=\upshape} % 取消所有列表默认距离
% 浮动环境设置
% 默认情况下, \LaTeX{} 要求每页的文字至少占据 20%,否则该页就只单独放置一个浮动环境,
% 而这通常不是我们想要的, 我们将这个要求降低到 5%.
\renewcommand*{\textfraction}{0.05}
% 有时如果多个浮动环境连续放在一起,
% 会将它们分在几个不同页,即使它们可在同一页放
% 得下. 我们可以通过修改 |\topfraction| 和 |\bottomfraction| 分别设置顶端和底端的浮
% 动环境的最大比例.
\renewcommand*{\topfraction}{0.9}
\renewcommand*{\bottomfraction}{0.8}
% 有时\LaTeX{}会把一个浮动环境单独放在一页,
% 我们要求这个环境至少要占据 85% 才能单独放在一页.
% 注意:  |\floatpagefraction| 的数值必须小于 |\topfraction|.
\renewcommand*{\floatpagefraction}{0.85}
% 关于图片 graphicx
% 如果图片没有指定后缀, 依次按下列顺序搜索
\DeclareGraphicsExtensions{.pdf,.eps,.jpg,.png}
% 设置图表搜索路径, 可以给图表文件夹取如下名字
\graphicspath{{figures/}{figure/}{pictures/}{picture/}{pic/}{pics/}{image/}{images/}}
\usepackage{mtpro2}
\usepackage[physics]{stys/physicx}
\usepackage{stys/Symbols}
\usepackage{amsfonts}
\DeclareFontFamily{U}{nxlmi}{}
\DeclareFontSubstitution{U}{nxlmi}{m}{it}
\DeclareFontShape{U}{nxlmi}{m}{it}{
<-6.3>    nxlmi05
<6.3-8.6> nxlmi07
<8.6->    nxlmi0
}{}

\DeclareFontShape{U}{nxlmi}{b}{it}{
<-6.3>    nxlbmi05
<6.3-8.6> nxlbmi07
<8.6->    nxlbmi0
}{}

\renewcommand{\partial}{{\text{\usefont{U}{nxlmi}{m}{it}\symbol{64}}\mspace{1mu}}}

%% 定义第一种定理
\mynewtheorem{
    defi={\textbf{定义}}[section]{interior style={left color=ReD!8,right color=ReD!5!CyaN!50}, borderline west={1.5mm}{0mm}{ReD}},
    thm={\textbf{定理}}[section]{interior style={left color=CyaN!80!black!20,right color=CyaN!80!black!15!CyaN!50}, borderline west={1.5mm}{0mm}{CyaN!80!black}},
    lem={\textbf{引理}}[section]{interior style={left color=BluE!8,right color=BluE!5!CyaN!50}, borderline west={1.5mm}{0mm}{BluE}},
    prop={\textbf{命题}}[section]{interior style={left color=OrangE!8,right color=OrangE!5!CyaN!50}, borderline west={1.5mm}{0mm}{OrangE}},
    exam={\textbf{题}}[chapter]{interior style={left color=DarkGreen!8,right color=DarkGreen!5!CyaN!50}, borderline west={1.5mm}{0mm}{DarkGreen}},
    cor={\textbf{推论}}[chapter]{interior style={left color=violet!8,right color=violet!5!CyaN!50}, borderline west={1.5mm}{0mm}{violet}},
}
\newtheorem*{remark}{\textbf{注}}
%% 定义第二种定理
% overlay unbroken=\my@theorem@overlay@unbroken{\theorem@name\ \thetcbthm}{额外的选项}
% overlay first=\my@theorem@overlay@first{\theorem@name\ \thetcbthm}{额外的选项}
%% 用户接口区
\definecolor{examback}{HTML}{e3e6e8}
\makeatletter
\mynewtcbtheorem{
    % 这个 theorem 是环境名
    theorem={
        counter=tcbthm, 
        the counter=\thesection.\arabic{tcbthm}, 
        name=定理, % 它保存到 \theorem@name 里
        thmcolor=purple5,
        autoref name=\bfseries 定理, 
        style={
        arc=3pt,breakable,enhanced,interior style={top color=purple5!5 ,middle color=purple5!1!nuanbai, bottom color=nuanbai},boxrule=0pt,top=8mm,
        fuzzy shadow={-0.6mm}{0.6mm}{0mm}{0.3mm}{white!50!gray},% 上
        fuzzy shadow={0.6mm}{-0.6mm}{0mm}{0.3mm}{fill=white!40!gray},%下
        opacityframe=0, opacityback=0.98,
        fontupper=\itshape, step={tcbthm},
        before pre=\smallskip, after app=\smallskip,
        overlay unbroken=\my@theorem@overlay@unbroken{\theorem@name\ \thetcbthm}{\theorem@thmcolor},
        overlay first=\my@theorem@overlay@first{\theorem@name\ \thetcbthm}{\theorem@thmcolor},
        overlay last=\my@theorem@overlay@last{\theorem@thmcolor},
        }
    },
    proposition={
        counter=tcbprop, 
        the counter=\thesection.\arabic{tcbprop}, 
        autoref name=\bfseries 命题, 
        style={
        arc=3pt,breakable,enhanced,interior style={top color=DeepPink!5 ,middle color=DeepPink!1!nuanbai, bottom color=nuanbai},boxrule=0pt,top=8mm,
        fuzzy shadow={-0.6mm}{0.6mm}{0mm}{0.3mm}{white!50!gray},% 上
        fuzzy shadow={0.6mm}{-0.6mm}{0mm}{0.3mm}{fill=white!40!gray},%下
        opacityframe=0, opacityback=0.98,
        fontupper=\itshape, step={tcbprop},
        before pre=\smallskip, after app=\smallskip,
        overlay unbroken=\my@theorem@overlay@unbroken{命题\ \thetcbprop}{DeepPink},
        overlay first=\my@theorem@overlay@first{命题\ \thetcbprop}{DeepPink},
        overlay last=\my@theorem@overlay@last{DeepPink},
        }
    },
    definition={
        counter=tcbdefi, 
        the counter=\thesection.\arabic{tcbdefi}, 
        autoref name=\bfseries 定义, 
        style={
        arc=3pt,breakable,enhanced,interior style={top color=紫棠!5 ,middle color=紫棠!1!nuanbai, bottom color=nuanbai},boxrule=0pt,top=8mm,
        fuzzy shadow={-0.6mm}{0.6mm}{0mm}{0.3mm}{white!50!gray},% 上
        fuzzy shadow={0.6mm}{-0.6mm}{0mm}{0.3mm}{fill=white!40!gray},%下
        opacityframe=0, opacityback=0.98,
        fontupper=\itshape, step={tcbdefi},
        before pre=\smallskip, after app=\smallskip,
        overlay unbroken=\my@theorem@overlay@unbroken{定义\ \thetcbdefi}{紫棠},
        overlay first=\my@theorem@overlay@first{定义\ \thetcbdefi}{紫棠},
        overlay last=\my@theorem@overlay@last{紫棠},
        }
    },
    lemma={
        counter=tcblem,
        the counter=\thesection.\arabic{tcblem},
        name=引理, 
        lemcolor=靛蓝, 
        autoref name=\bfseries 引理,
        style={
        arc=0mm,breakable,enhanced,interior style={top color=靛蓝!5 ,middle color=靛蓝!1!nuanbai, bottom color=nuanbai},arc=3pt,boxrule=0pt,top=7mm,bottom=5mm,
        fuzzy shadow={-0.6mm}{0.6mm}{0mm}{0.3mm}{white!50!gray},% 上
        fuzzy shadow={0.6mm}{-0.6mm}{0mm}{0.3mm}{fill=white!40!gray},%下
        opacityframe=0, opacityback=0.98,
        fontupper=\normalsize,step={tcblem},
        before pre=\smallskip, after app=\smallskip,
        overlay unbroken=\my@lemma@overlay@unbroken{\lemma@name\ \thetcblem}{\lemma@lemcolor},
        overlay first=\my@lemma@overlay@first{\lemma@name\ \thetcblem}{\lemma@lemcolor},
        overlay last=\my@lemma@overlay@last{\lemma@lemcolor},
        }
    },
    corollary={
        counter=tcbcor,
        the counter=\thesection.\arabic{tcbcor},
        autoref name=\bfseries 推论,
        style={
        arc=0mm,breakable,enhanced,interior style={top color=茶色!5 ,middle color=茶色!1!nuanbai, bottom color=nuanbai},arc=3pt,boxrule=0pt,top=7mm,bottom=5mm,
        fuzzy shadow={-0.6mm}{0.6mm}{0mm}{0.3mm}{white!50!gray},% 上
        fuzzy shadow={0.6mm}{-0.6mm}{0mm}{0.3mm}{fill=white!40!gray},%下
        opacityframe=0, opacityback=0.98,
        fontupper=\normalsize,step={tcbcor},
        before pre=\smallskip, after app=\smallskip,
        overlay unbroken=\my@lemma@overlay@unbroken{推论\ \thetcbcor}{茶色},
        overlay first=\my@lemma@overlay@first{推论\ \thetcbcor}{茶色},
        overlay last=\my@lemma@overlay@last{茶色},
        }
    },
    example={
        counter=tcbexam,
        the counter=\thesection.\arabic{tcbexam},
        autoref name=\bfseries 例题,
        style={
        arc=0mm,breakable,enhanced,interior style={top color=黛绿!5 ,middle color=黛绿!1!nuanbai, bottom color=nuanbai},arc=3pt,boxrule=0pt,top=7mm,bottom=5mm,
        fuzzy shadow={-0.6mm}{0.6mm}{0mm}{0.3mm}{white!50!gray},% 上
        fuzzy shadow={0.6mm}{-0.6mm}{0mm}{0.3mm}{fill=white!40!gray},%下
        opacityframe=0, opacityback=0.98,
        fontupper=\normalsize,step={tcbexam},
        before pre=\smallskip, after app=\smallskip,
        overlay unbroken=\my@lemma@overlay@unbroken{例题\ \thetcbexam}{黛绿},
        overlay first=\my@lemma@overlay@first{例题\ \thetcbexam}{黛绿},
        overlay last=\my@lemma@overlay@last{黛绿},
        }
    },
}
\makeatother
\newenvironment{note}[1][\bf 笔记:]{\Line\uuline{#1} }{\Line}
\renewcommand{\Line}{\noindent\\\tikz\draw[line width=0.65pt,gray!80,dashed] (0,0)--++(.99\linewidth,0);\par}
\newenvironment{key}[1][]{\begin{fancybox}{#1}}{\end{fancybox}}
\newcommand{\Diff}[2][]{\frac{\partial #1}{\partial #2}}
\newcommand{\Dif}[2]{\frac{\dd #1}{\dd #2}}
\newcommand{\pr}{^\prime}
\usepackage{extarrows}
\usetikzlibrary{tikzmark}
% \arrowname{super-script}
% \arrowname[sub-script]{super-script}
\begin{document}
\thispagestyle{empty}
\title{Reviewing Material}
\subtitle{信息技术学考导引试题详解}
\edition{First Edition}
\bookseries{Preperation For The Final Test}
\author{xaero}
\pressname{}
\presslogo{figures/hglogo.pdf}
\coverimage{figures/pexels-alex-andrews-816608.jpg}
\makecover

\definecolor{bg}{HTML}{e0e0e0}
\definecolor{fg}{HTML}{455a64}
\colorlet{outermarginbgcolor}{bg}
\colorlet{outermarginfgcolor}{fg}
\colorlet{framegolden}{fg}
\colorlet{framegray}{黛绿!5}

%\thispagestyle{empty}
\begin{titlepage}
\thispagestyle{empty}
    \begin{center}
    
    %% Print the bookseries
    {\makeatletter\vspace*{5mm}
    \ifdefvoid{\@bookseries}{}{\bigskip\normalfont\fontsize{20}{20}\selectfont\@bookseries}
    \makeatother}

        \bigskip
        \bigskip

    %% Print the title
    {\makeatletter
    \fontsize{18}{20}\rmfamily\selectfont\@title \\\vspace*{5mm}
    \fontsize{32}{35}\rmfamily\bfseries\selectfont 2023信息技术学考导引试题详解
    \makeatother}
    
    \bigskip
    \bigskip
    
    \bigskip
    \bigskip
    \bigskip

    %% Print the name of the author
    {\makeatletter
    \fontsize{16}{25}\rmfamily\selectfont 陈\hspace{5mm}功\hspace{5mm}杰
    \makeatother}
    
    \bigskip
    \bigskip

    
    \vfill
    
    
    \bigskip
    \bigskip
    
    {\makeatletter
    %\fontsize{25}{25}\rmfamily\selectfont\@pressname
    \vspace*{-5mm}
    \centering
    \includegraphics[width=0.4\linewidth]{figures/hglogoblack}    
    \makeatother}
    \end{center}
    
    \end{titlepage}
    \let\cleardoublepage\clearpage

    \thispagestyle{empty}
    \begin{center}
        {
        \vspace*{5cm}
        \fontsize{14}{20}\rmfamily\selectfont    源\hspace{1em}码\hspace{1em}获\hspace{1em}取}\\ 
        \bigskip
        \bigskip
        \href{https://github.com/xaero7/HG_DaoYin}{https://github.com/xaero7/HG\_DaoYin}


    \end{center}







\frontmatter
\pagenumbering{Roman}
%\input{chapters/preface}
\thispagestyle{empty}
\tableofcontents\let\cleardoublepage\clearpage

\mainmatter
\pagenumbering{arabic}
\partimage{figures/5.png}
%\partabstract{\hspace{2em}本书系统地论述了微分几何的基本知识. 作者用前3章,  以及第6章共计4章的篇幅介绍了流形、多重线性函数、向量场、外微分、李群和活动标架等基本知识和工具.  基于上述基础知识,  论述了微分几何的核心问题,  即联络、黎曼几何、以及曲面论. 第7章是当前十分活跃的研究领域——复流形. 陈省身先生是此研究领域的大家,  此章包含有作者独到、深刻的见解和简捷、有效的方法. 第8章的Finsler几何是本书第2版新增加的一章,  它是陈省身先生近年来一直倡导的研究课题,  其中Chern联络具有突出的性质,  它使得黎曼几何成为Finsler几何的特殊情形. 最后两个附录,  介绍了大范围曲线论和曲面论,  以及微分几何与理论物理关系的论述,  为这两个活跃的前沿领域提出了不少进一步的研究课题.

%\hspace*{2em} 本书的作者之一是已故数学家陈省身先生,  他开创并领导着整体微分几何、纤维丛微分几何、“陈省身示性类”等领域的研究,  他是第一个获得世界数学界最高荣誉“沃尔夫奖”的华人,  被称为“当今最伟大的数学家”,  被国际数学界尊为“微分几何之父”. }

%% \part{微分几何讲义一陈省身}

%\chapter{必修1 $\cdot$ 数据与计算}\index{数据与计算}
\section{数据与信息}\index{数据与信息}

\begin{enumerate}
%% ============= 1
\item 答案：D。考查数据与信息概念的理解，数据、数字的差别。
	\begin{enumerate}[label=\Alph*.]
	\item 数据是对客观事物的符号表示，如图形符号、数字、字母等。在计算机中的表示形式可以是文字、图形、图像、音频、视频等。
	\item 数据可以加过处理，但显然你可以让他失去原有的价值。
	\item 数字放到特定的环境、语境下才有意义，即要有上下文才有含义。
	\item 正确。
	\end{enumerate}

%% ============= 2
\item 答案：B。考查信息概念的理解，信息的特征。
	\begin{enumerate}[label=\Alph*.]
	\item 实验误差是测量值和真实值之间的偏差，不是虚假信息。
	\item 正确。
	\item 同一个信息对于不同的人价值可能不一样。
	\item 信息是信号、消息中所包含的含义，必须依附与数字、文字、图形、图像等载体。
	\end{enumerate}

%% ============= 3
\item 答案：D。考查信息概念的理解，信息的特征。
	\begin{enumerate}[label=\Alph*.]
	\item 互联网上只有已数字化的信息，没有数字化当然查不到。
	\item 知识的获得是人利用自身已有的知识对信息进行加工，进而将新的信息纳入自己的知识结构的过程。检索到也只是看到，并不一定已内化成自己的知识。
	\item 天才也要记单词啊。
	\item 正确。
	\end{enumerate}

%% ============= 4
\item 答案：A。
	\begin{enumerate}[label=\Alph*.]
	\item 。
	\item 。
	\item 。
	\item 。
	\end{enumerate}

%% ============= 5
\item 答案：C。
	\begin{enumerate}[label=\Alph*.]
	\item 若化成十进制计算：$10H=16D, 10B = 2D, 16D+2D=18D$。
	\item $1AH + 2AH = 44H$，注意十六进制下$A + A$等于$14$。
	\item 正确。
	\item 若化成十六进制计算：$10D + 10B = AH + 2H = CH$，即十六进制的值是$C$。
	\end{enumerate}

%% ============= 6
\item 答案：D。
	\begin{enumerate}[label=\Alph*.]
	\item 。
	\item 。
	\item 。
	\item 。
	\end{enumerate}

%% ============= 7
\item 答案：D。
	\begin{enumerate}[label=\Alph*.]
	\item 。
	\item 。
	\item 。
	\item 。
	\end{enumerate}

%% ============= 8
\item 答案：D。
	\begin{enumerate}[label=\Alph*.]
	\item 。
	\item 。
	\item 。
	\item 。
	\end{enumerate}

%% ============= 9
\item 答案：B。
	\begin{enumerate}[label=\Alph*.]
	\item 。
	\item 。
	\item 。
	\item 。
	\end{enumerate}

%% ============= 10
\item 答案：B。
	\begin{enumerate}[label=\Alph*.]
	\item 。
	\item 。
	\item 。
	\item 。
	\end{enumerate}

%% ============= 11
\item 考查信息编码、容量计算。
	\begin{enumerate}[label=$(\arabic*)$]
	\item {\kaishu 视频容量=每帧图像容量$\times$帧频，每帧图像容量=像素点数$\times$量化位数}。依题意，单张图像容量是：$\dfrac{1280 \times 720 \times 24}{8 \times 1024 \times 1024} \approx 2.64$MB。因此视频容量是：$2.64 \times 5 \times 60 \times 25 = 19800$MB。压缩比至少是$39.6:1$才能压缩到500MB以内。答案是$40:1$。
	\item 加入数据不会改变原来的压缩比，相当于不会改变原先的压缩编码方式（真因为如此，加入音频后的视频容量增加，压缩比不变的话，压缩之后的作品容量也增加，势必会超过500MB，因此必须重新设定新的、更大的压缩比才能压缩到500MB以内，这应该是本题想考查的一个实际应用情景）。
	\item 压缩会使画面不清晰，原因压缩比太大，或者压缩算法太差。改进办法是可以换一种压缩算法（换一个压缩软件），或者保证内容完整的前提下，减少画面尺寸、缩短时长等。
	\end{enumerate}

\end{enumerate}


\newpage

%\section{算法与程序设计}\index{算法与程序设计}

\begin{enumerate}
%% ============= 1
\item 答案：A。

%% ============= 2
\item 答案：A。

%% ============= 3
\item 答案：D。

%% ============= 4
\item 答案：C。

%% ============= 5
\item 答案：B。

%% ============= 6
\item 答案：C。选项A、B都可以用解析法，用公式计算。

%% ============= 7
\item 答案：D。使用列表法模拟流程图的执行。

%% ============= 8
\item 答案：A。

%% ============= 9
\item 答案：D。

%% ============= 10
\item 答案：C。计算机表达式：写在同一行，没有上、小标的数学写法，一般只有字母、数字等标识符，$\pi$不是英文字母。

%% ============= 11
\item 答案：B。

%% ============= 12
\item 答案：B。

%% ============= 13
\item 答案：A。考查Python循环语句、双重循环程序的阅读理解。
\columnratio{0.46}
\begin{paracol}{2}
\begin{lstlisting}[numbers=left]
for i in range(1, 7):
    for j in range(1, 7):
        if j <= i:
            print(j, end=" ")
        else:
            print("", end="")
    print()
\end{lstlisting}
\switchcolumn
固定第$1$行处的外层循环$i$的值为$1$时，内层循环$j$从$1$变化到$6$，对于每一个$j$，当$j \le i$时输出$j$的值，否则输出空值。因此当$i=1$时，输出$1$，然后换行；当$i=2$时，输出$1 \; 2$然后换行，当$i=3$时输出$1 \; 2 \; 3$然后换行……，答案选A。
\end{paracol}

%% ============= 14
\item 答案：（1）确定性，让输入错误次数加一。（2）②

%% ============= 15
\item 答案：（1）$|s * s - 2| \ge 0.0001?$，算法中可以用数学符号，也可以用计算机表达式。（2）B。

%% ============= 16
\item 答案：（1）①②⑤ （2）①顺序结构 ②$i=i+m$

%% ============= 17
\item 答案：① w <= 20 ② 0.6*15+5*1.4+(w-20)*2.1

%% ============= 18
\item 答案：① a=b ②str(c)或str(b)

%% ============= 19
\item 答案：① n=len(h) ② code[a]*16**(n-i-1)

%% ============= 20
\item 考查应用Python程序解决实际问题的能力。考查字符串的处理与应用。
	\begin{enumerate}[label=$(\arabic*)$]
	\item 考查题意的理解，这是理解题目情景的关键。“we put the bed in the bedroom”中有两处“bed”，会被替换两次。	
	\item 阅读与推导过程：
\setcounter{qnumber}{1}
\begin{lstlisting}[numbers=left]
text = input("输入原文字符串: ")
key = input("输入要查找的字符串: ")
rs = input("输入替换字符串:")
result = "";  count = 0;  i = 0;  n = len(text)
while i < n - len(key) + 1:
 s = text[`\clozeblank{2}`]
 if key == s:
     result += rs
     count += 1
     i += len(key)
 else:
     result += text[i]
     i += 1
`\clozeblank{2}`
if count > 0:
 print("替换的次数: ", count)
 print("替换后的结果: ", result)
else:
 print("要查找的内容不存在")
\end{lstlisting}
		\begin{enumerate}[label=$(\alph*)$]
		\item 第5行的循环和$n$有关，而$n$是原文的长度，因此第5行的循环是在扫描原文的每个字符。
		\item 从$i$的变化上看，当第8行两个字符相等时，$i$往后移动与key一样的长度；当两个字符不等时，$i$往后移动1个字符长度，所以$i$是指示了原文text中待比较字符串的索引位置信息。
		\item 循环中第7行判定了key是否与$s$相等，那么$s$就需要从原文中截取一个字符串，再与key作相等比较，因此第①空应该填写原文字符串的切片，切片的起始值是当前$i$的值，切片的长度应该与key的长度相等，于是第①空答案是\lstinline|i:i+len(key)|。
		\item 在解题时一定要用样例带入后阅读，比如原文\lstinline|text="Abedrbedom"|，待替换字符串\lstinline|key="bed"|，那么当$i=1,5$时分别找到两处“bed”，如下图所示。当$i$指向8号位置时，剩余字符串已不足3位（即待查找值key的长度），也就无需继续循环，这也是第5行while循环条件是\lstinline|i < n - len(key) + 1|而不是\lstinline|i < n|的原因。但是这样带来的后果是剩余的字符串无法原样连接到result中取，如下图中的最后两个字符“om”。因此需要在循环结束时，第14处将剩余字符串连接到最终结果上。第②空的答案是\lstinline|result += text[i:]|，其中切片的终止端点写明$n$亦可。
\begin{lstlisting}
0 1 2 3 4 5 6 7 8 9
A b e d r b e d o m
  ↑
  i=1时找到第一处，字符串替换后，i=i+3，指向4号位置
----------------------------------------------------------------
0 1 2 3 4 5 6 7 8 9
A b e d r b e d o m
          ↑
          i=5时找到第二处，字符串替换后，i=i+3，指向8号位置
\end{lstlisting}
		\end{enumerate}
	\end{enumerate}

%% ============= 21
\item 考查应用Python程序解决实际问题的能力。考查字符串的处理与应用。
	\begin{enumerate}[label=$(\arabic*)$]
	\item “Good”中的四个字母来自键盘不同的两行按键，故答案是No。
	\item 阅读与推导过程：
\setcounter{qnumber}{1}
\begin{lstlisting}[numbers=left]
def to_lower(ch):       # 转小写字母
    if ch >= "A" and ch <= "Z":
        return chr(ord(ch) + 32)
    else:
        `\clozeblank{2}`
line_1 = "qwertyuiop"   # 键盘第一行字母
line_2 = "asdfghjkl"    # 键盘第二行字母
line_3 = "zxcvbnm"      # 键盘第三行字母
char = input()
c1 = 0
c2 = 0
c3 = 0
n = len(char)
for ch in char:
    ch = `\clozeblank{2}`
    if ch in line_1:
        c1 += 1
    elif ch in line_2:
        c2 += 2
    elif ch in line_3:
        c3 += 3
if c1 == n or c2 == n * 2 or c3 == n * 3:
    print("yes")
else:
    print("no")
\end{lstlisting}
	\begin{enumerate}[label=$(\alph*)$]
	\item 第①空容易填：第2行if语句判定了大写字母，第3行将其转成小写字母，因此非大写字母时直接返回ch，答案是\lstinline|return ch|。
	\item 第14行循环语句遍历提取了输入字符串char中的每个字符，在第16、18、20行分别判定了是否是键盘上哪一行的字母：第一行则$c1$加1，第二行则$c2$加2，第三行则$c3$加3，由此可以断定，在第②空处需要将字母规范化——统一转成小写字母，这就需要调用\lstinline|to_lower()|函数，因此答案是\lstinline|to_lower(ch)|，参数是当前扫描到的字符ch。
	\item 对于第22行条件的理解：如果输入字符char都来自键盘第一行，那么$c1$的值与char字符串长度相等，因此$c1$每次都加1；同理，如果都来自第二行，则$c2$的值是char字符串长度的两倍，因为$c2$每次都加2；$c3$的值亦同理。
	\end{enumerate}
	\item \lstinline|c1 += 1|的含义是第一行的字符数量增加1
	\end{enumerate}

%% ============= 22
\item 考查应用Python程序解决实际问题的能力。考查随机数函数、枚举算法。
	\begin{enumerate}[label=$(\alph*)$]
\setcounter{qnumber}{1}
\begin{lstlisting}[numbers=left]
import random
n = int(input("请输入要产生的英文字符串长度: "))
s = ""
for i in range(n):
	# randint(1,58):随机生成一个[1,58]范围内整数，字母A的ASCII码值为65
    s += chr(64 + random.randint(1, 58))  
print(s)
ans = input("请按样例输入: ")
c = 0
for i in range(n):
    if `\clozeblank{2}`:
        c += 1
p = `\clozeblank{2}`
print("正确数量: ", c, "，正确率为: ", p, "%")
\end{lstlisting}
	\item 判断两个字符串有多少个字符相同，可以用枚举算法：遍历每个字符串的每一位，分别判定是否相等。
	\item 由第10行的for循环语句知，$i$取遍了$[0,-1]$的每个数，这相当于字符串的索引值。而第11行处的条件成立时，变量$c$的值加1，又由第14行的输出可知$c$是正确单词的个数。因此第①空是判定两个字母是否相等，原始字符串是s，用户输入字符串是ans，因此答案是\lstinline|s[i] == ans[i]|。
	\item 变量$p$是什么？同样可以看14行的输出语句——$p$是正确率百分比。因此$p$的计算方式是正确个数除上总个数，答案可以是\lstinline|c / n * 100|。题意没有说如何保留小数，也没有输出示例，因此这个答案也可以。参考答案是\lstinline|int(c / n * 100 + 0.5)|，它的功能是四舍五入保留整数。
	\end{enumerate}


%% ============= 23
\item 考查应用Python程序解决实际问题的能力。考查进制转换解析算法、字符串应用。
	\begin{enumerate}[label=$(\alph*)$]
\setcounter{qnumber}{1}
\begin{lstlisting}[numbers=left]
def conv(s):
    ans = ""
    if s > "9":
        `\clozeblank{2}`
    else:
        s = int(s)
    while s > 0:
        k = s % 2
        s //= 2
        `\clozeblank{2}`
    for i in range(4 - len(ans)):
        ans = "0" + ans
    return ans
s = "2A08:CCD6:0088:108A:0011:0002:202F:AA05"
ans = ""
flag = False
for i in s:
    if i == ":":
        `\clozeblank{2}`
        ans += i
    elif i != "0" or flag == True :
        `\clozeblank{2}`
        flag = True
print("原IPv6地址为:", s)
print("去前导零后:", ans)
\end{lstlisting}
	\item 当自定义函数比较复杂时，可以从主程序开始阅读。那么从第14行开始阅读程序：第17行遍历取出了字符串$s$中的每个字符，$s$是个十进制模式的IPv6地址字符串，依题意需要将它转成二进制模式，可以猜测本题处理思路就是逐个取出$s$的每个字符并转换成二进制并输出结果。
	\item 第18行是$i$的值是冒号的情况，这意味着冒号前面一段IPv6已转换完，$i$中的值直接连接到最后结果字符串ans变量的后面，ans变量的功能也还是从最后一行的输出语句中得到。但是第③空还不知道填什么。
	\item 第21行是说当$i$不是0，或者flag的值是True时执行第22、23行程序。容易相当，当$i$非0时必然要转成二进制格式，因此这里需要调用前面的conv()函数，调用结果应该是$i$字符对应的二进制串，同样要把二进制串连接到ans变量后面，于是第④空的答案应该是\lstinline|ans += conv(i)|，其中$i$就是当前需要转换成二进制的十六进制字符。
	\item 从第21行的elif判定结果看，当$i \neq 0$时，flag的结果会变成True，一个隐含的情况是当$i=0$但flag值是True的时候，也会执行21、22两行代码，即也会将该“0”转成二进制串。结合题意“前导零可以省略”可知，非前导（中间的）零需要转换。由此断定flag的值为True表示当前有非前导的零（需要转换）；flag的值为False时，若出现零则是前导的零。再结合第16行f初值False可知，这样的推论是合适的。
	\item 因此，第③空是出现冒号后，下一次得到的字符串若是零，该零必然是前导的零，于是这里填\lstinline|flag = False|。
	\item 转到conv()函数，可以断定参数变量$s$是待转换的十六进制字符串，这在第8、9两行的循环模2取余也可以得到验证（转二进制的方法就是除二取余法）。
	\item 但是$s$是十六进制，除二取余之前需要转成十进制。阅读第$1 \sim 6$行的if语句可以看出，若$s \le 9$，则直接取整（此时十六进制值与十进制相等）；否则要把“A”转成10，把“B”转成11……把“F”转成15。把字母转整数可以用ASCII码函数ord()，本题答案是\lstinline|ord(s) - 65 + 10|。
	\item 第②空是将余数$k$连接成二进制串的语句，注意最先除二取余得到的余数是最低位，最后得到的余数是最高位，因此本题答案需要将$k$转成字符串后连接到答案变量ans的前面。本题填\lstinline|ans = str(k) + ans|。	
	\end{enumerate}


%% ============= 24
\item 考查应用Python程序解决实际问题的能力，考查列表的应用。
	\begin{enumerate}[label=$(\alph*)$]
\setcounter{qnumber}{1}
\begin{lstlisting}[numbers=left]
import random as rd
data = [180,283,385,170,276,384,180,285,380,190,295,390,170,272,372]
s = [0, 0, 0]                   # 存储3个作品的得分
ans = []                        # 存储并列最高平均分的作品号
maxb = 0
for i in range(len(data)):
    zp = `\clozeblank{2}`   # 分离出作品编号
    fs = data[i] % 100
    `\clozeblank{2}`        # 累加当前作品的得分
for j in range(3):
    `\clozeblank{2}`
    print("作品", j+1, "平均分为", s[j])
    if s[j] > maxb:
        maxb = s[j]
for z in range(3):              # 查找并列最高平均分
    if `\clozeblank{2}`:
        ans.append(z + 1)       # 将数据添加到列表ans尾部
print("平均分最高作品号是: ", ans)
\end{lstlisting}
	\item 由第6行的循环范围知，该for循环遍历了data列表的每个元素，$i$是其索引值。
	\item 分离字符串可以用切片，分离整数的各个数位可以用对十取余数，或者整除十。从循环中的第8行知，\lstinline|data[i] % 100|就是该整数的十位和各位上的数，即第$i$号数据中的得分值。第②空，分离百位数可以模仿这写\lstinline|data[i] // 100|，除以100之后的整数部分就是作品号。
	\item 按注释，第②空应是累加fs的值到列表$s$中的某个位置上，这个位置应该与作品编号有关。注意到第7行分离出来的作品编号都是从“1”开始数的，而$s$的索引值是从“0”开始计的，因此需要修正，答案是\lstinline|s[zp-1] += fs|。
	\item 因为有3个作品，因此第10行的3次循环应该遍历了每个作品，并输出了相应的平均分。也就是说，第12行输出的s[j]是第$j$号作品的平均分，而之前$s$中保存的是每个作品的总分，因此第③空需要求平均分，答案是\lstinline|s[j] /= 5|。
	\item 第13行是打擂算法：若当前平均分s[j]大于“擂台”上的数maxb，则让s[j]留在擂台上。所以maxb保存的是最大值。
	\item 第④空的条件成立时将$z$存入列表ans，而由最后一行的输出结合输出图示看，列表ans中存放了所有得分都是最高的作品编号。第④空就得填写某个作品均值与maxb相等时执行插入操作，答案是\lstinline|s[z] == maxb|。
	\end{enumerate}






































\end{enumerate}


\newpage

\section{数据处理与应用}\index{数据处理与应用}

\begin{enumerate}
%% ============= 1
\item 答案：D。考查数据整理方法与目的。

%% ============= 2
\item 答案：C。
	\begin{enumerate}[label=\Alph*.]
	\item 实验误差是测量值和真实值之间的偏差，不是虚假信息。
	\item 正确。
	\item 同一个信息对于不同的人价值可能不一样。
	\item 信息是信号、消息中所包含的含义，必须依附与数字、文字、图形、图像等载体。
	\end{enumerate}

%% ============= 3
\item 答案：B。
	\begin{enumerate}[label=\Alph*.]
	\item 互联网上只有已数字化的信息，没有数字化当然查不到。
	\item 知识的获得是人利用自身已有的知识对信息进行加工，进而将新的信息纳入自己的知识结构的过程。检索到也只是看到，并不一定已内化成自己的知识。
	\item 天才也要记单词啊。
	\item 正确。
	\end{enumerate}

%% ============= 4
\item 答案：B。考查pandas数据处理drop()函数、groupby()函数功能的理解。注释如下：
\begin{lstlisting}[numbers=left]
import pandas as pd									# 导入并使用pd作为别名
df = pd.read_csv("mnxk.csv", sep=",")				# 读取数据
df1 = df.drop("已选科目数", axis=1)				    # 删除“已选科目数”列
print(df.head())									# 打印df的前5行
print(df1.head())									# 打印df1的前5行
sc=df1.groupby("班级", as_index=False).count()	    # 按“班级”分组
m = len(df)											# df的行数
n = len(df1)										# df1的行数
print(sc)											# 分组后的数据
\end{lstlisting}
注意pandas的很多操作处理后原始数据都不会改变。比如，第3行drop()函数删除了“已选科目数”列，参数“axis=1”指明了这是列而不是行。该函数调用后，产生了一个新的数据集合并赋值给对象df1，而原始的数据集合df未曾变化。选项B就考查了pandas数据处理的这个特点：第4行打印的结果是原始数据的前5行，包含“已选科目数”这列数据，而第5行的输出的5行数据虽然大部分与前面相同，但不含“已选科目数”这列数据。选项C考查的是df和df1数据对象的行数是否相同，由于没有删除行，行数必然是一样的。选项D，第6行的分组可以让相同班级的数据合并成一行数据，这个数据的每列数据是原先该列数据的非空单元格个个数（即count()函数的功能）。如，若原始数据如左侧所示，则执行第6行分组语句后的结果如右侧所示。在右侧数据中，“1班”的“Name”值是4，表示原始数据中1班“Name”列数据非空单元格个有4个；“1班”的“物理”值是2，表示原始数据中1班“物理”列数据非空单元格个有2个（相当于1班有两个2人选了物理）。
\begin{lstlisting}
   班级 Name  物理 历史 技术 化学          班级 Name 物理 历史 技术 化学
0  1班  张三丰   1   1   1                0  1班  4   2   1   2   2
1  2班   郭靖    1       1                1  2班  2   1   0   1   0
2  1班  小龙女   1           1            2  3班  2   2   1   1   2
3  2班  李秋水                
4  3班   杨过    1       1   1
5  1班  令狐冲                
6  3班  任我行   1   1       1
7  1班   黄蓉            1   1
\end{lstlisting}

%% ============= 5
\item 答案：D。
	\begin{enumerate}[label=\Alph*.]
	\item 若化成十进制计算：$10H=16D, 10B = 2D, 16D+2D=18D$。
	\item $1AH + 2AH = 44H$，注意十六进制下$A + A$等于$14$。
	\item 正确。
	\item 若化成十六进制计算：$10D + 10B = AH + 2H = CH$，即十六进制的值是$C$。
	\end{enumerate}

%% ============= 6
\item 答案：B。
	\begin{enumerate}[label=\Alph*.]
	\item 。
	\item 。
	\item 。
	\item 。
	\end{enumerate}

%% ============= 7
\item 答案：B。
	\begin{enumerate}[label=\Alph*.]
	\item 。
	\item 。
	\item 。
	\item 。
	\end{enumerate}

%% ============= 8
\item 答案：D。
	\begin{enumerate}[label=\Alph*.]
	\item 。
	\item 。
	\item 。
	\item 。
	\end{enumerate}

%% ============= 9
\item 答案：C。
	\begin{enumerate}[label=\Alph*.]
	\item 。
	\item 。
	\item 。
	\item 。
	\end{enumerate}

%% ============= 10
\item 考查pandas数据处理与应用。
	\begin{enumerate}[label=$(\arabic*)$]
	\item 考查数据处理的实际用途，帮助理解题目情景。
	\item 考查pandas数据格式的识别。
\setcounter{qnumber}{1}
\begin{lstlisting}[numbers=left]
def s_review(c):    
    for r in range(df.shape[0]):            # 批阅1个单选题
        if df.at[r, qnum[c]] == `\clozeblank{2}`:
            tmp = 3
            df.at[r, score[c-2]] = tmp
            df.at[r, score[10]] += tmp      # 计算总分，存入"sum"列
qnum = df.columns
sans = "BDCABDDBCB"                         # 本次作业的标准答案
score=["sc1","sc2","sc3","sc4","sc5","sc6","sc7","sc8","sc9","sc10","sum"]
for c in score:
    df[c] = 0
for c in range(2,12):                       # 逐题批阅
    `\clozeblank{2}`
print(df)
df.to_excel("客观题成绩.xlsx", index=False)  # 保存结果
\end{lstlisting}
	\item 解题过程：
		\begin{enumerate}[label=$(\alph*)$]
		\item 从第7行主程序开始阅读，对pandas程序阅读，一定要直到变量保存的数据是什么？数据的结构是怎样的？
		\item 第7行qnum保存了数据对象df的所有列名称，即\lstinline|["name", "snum", "ans1", "ans2", ..., "ans10"]|。第10行的循环，结合第9行score列表中的数据可以知道第11行在数据对象df中新增了很多数据列，列名称分别是“sc1”、“sc2”、“sc3”、……、“sum”，每列的值都是0。这也是pandas的特点，数据列直接可以参与算术运算、关系运算和赋值操作，每种操作都可以将该列的所有行都进行相应处理。
		\item 第②空处所进行的循环是逐题批阅。原始表格数据中一题就是一列数据，列序号是$2 \sim 12$，刚好能对上这里的循环范围。因此，第12行的循环变量$c$相当于列序号——不过，pandas需要的是列名称，这就需要qnum中对应的列名称来引用原始数据了。这里需要调用第一行的s\_review()函数。
		\item 阅读s\_review()函数。第2行df.shape可以返回数据对象df的维度“形状”：行数（df.shape[0]）和列数（df.shape[1]），因此$r$就是行索引号。由\lstinline|df.at[r, qnum[c]]|操作可知qnum[c]必然是列名称，结合前面的分析可以知道$c$必然是列序号。由于列表qnum中索引$2$号的列名称才是第一题名称“ans1”，因此$c$期望的值也应该从$2$开始。那么第②空的函数调用就好办了：函数名已知的，参数作用也推知了，所以答案应该是\lstinline|s_revieww(c)|，就让$c$的值从2开始传递、调用函数。另外从s\_review()函数的结构上看，它有return语句返回值，所以这空也无需考虑赋值——直接调用即可。
		\item 再回到第3行程序，\lstinline|df.at[r, qnum[c]]|取得了$c$列每个人填写的答案，它们需要与标准答案做比较，标准答案保存再sans字符串中，它的索引号是从0开始的，所以第①空的答案是\lstinline|sans[c-2]|。
		\item 第5行的程序是将该行（第$r$行）对应的得分列赋值为tmp分分值（如“ans1”列对应的分值是“sc1”列）。第6行的程序是将该分值累加到它的总分中去（即“sum”列，它的值是10个选择题的得分累加而来）。
		\end{enumerate}
	\end{enumerate}


%% ============= 11
\item 考查文本数据处理、分词、字符串统计与字典的应用。
	\begin{enumerate}[label=$(\arabic*)$]
\setcounter{qnumber}{1}
\begin{lstlisting}[numbers=left]
import jieba                                     # 导入jieba模块
import pandas as pd
text = open("news.txt", encoding="utf-8").read() # 打开文本文件
words = jieba.lcut(text, cut_all=False)          # 分词
counts = {}
for name in words:
    if len(name) != 1 and not ("a" < name[0] < "z") and not ("0" < name[0] < "9"):
        if name in counts:
            counts[name] += 1                    # 词语已出现过
        else:
            counts[name] = 1                     # 词语第一次出现
# 字典转化为DataFrame格式存储
df = pd.DataFrame(list(counts.items()), columns=["词", "次数"])
df = df.sort_values("次数", ascending=False)     # 按“次数”降序排序
print(df)
\end{lstlisting}
	\item jieba是目前常用的分词模块，它是一个基于词典分词的模块。模块导入后，程序再第3行通过python的内置open()函数打开了文本文件，read()函数可以读取文件中的所有数据。第4行调用了jieba的lcut()函数进行分词，函数名中的“l”表示分词结果数据是一个列表（即list，这里了解即可，无需记忆），函数的cut\_all参数设定为False表示是精准分词，不会分隔“词中词”，当该参数设定为True时表示全模式分词，会分隔所有词。如“中华人民共和国”，False模式下结果是一个词\lstinline|["中华人民共和国"]|，True模式下会有多个词\lstinline|["中华", "人民", "共和国", "中华人民共和国"]|。对于words列表中的每个词，第7行程序过滤掉了单字、字母开头的、数字开头的字符串，因此答案选C。
	\item 本小题考查jieba分词的规则特点，因为它是用现有的词典进行分词的，因此想要添加一个新词时，只需在分词前添加该词再进行分词即可。具体可以通过\lstinline|jieba.add_word( "公益活动")|来添加该词。
	\item 第8行程序先判定某个单词name是否在字典counts的键名中出现过，如果出现过，则直接根据该键名取出其键值，然后加1后仍然存放在该键名上。else分支就是该键名第一次出现，因此该键值初始为1，答案是\lstinline|counts[name] = 1|。
	\end{enumerate}


%% ============= 12
\item 考查文本数据处理、分词、字符串统计与字典的应用。
	\begin{enumerate}[label=$(\arabic*)$]
\setcounter{qnumber}{1}
\begin{lstlisting}[numbers=left]

\end{lstlisting}
	\item 
	\end{enumerate}











\end{enumerate}


\newpage







\end{document} 